% ======================================================================
% Section 5 — The group-ring translation
% ======================================================================

\section{The group-ring translation}\label{sec:groupring}

The covering lemmas of the $p^2$ programme (Lemmas A and B of the companion
notes) are statements about three arithmetic progressions covering $\Z_p$.
This section re-encodes them \emph{exactly} as a statement about one sparse
identity in the group ring $\Z[\Z_p]\cong\Z[z]/(z^p-1)$: the obstruction
becomes the coefficient-sign pattern of a quotient of a $24$-term element by
an $8$-term element. The translation is proved (Theorem~\ref{thm:GT}), not
merely verified, and it is unconditional --- the identity holds for
\emph{every} configuration, and covering is exactly a positivity condition on
its solution. The machine verification then establishes, over an independent
code path, a census law: at every prime tested, the critical coverings fall
into exactly four difference families with placement counts \emph{independent
of $p$} (\S\ref{sec:gt-census}). This is the same pinning phenomenon that the
shear-rigidity argument measured fiberwise, now visible as a statement about
the identity itself, and it converts the programme's computational inputs
from vague debt into precisely stated, named lemmas with identified
discharge paths (\S\ref{sec:gt-support}) --- one of which, the foot-count
lemma, is now proved (Lemma~\ref{lem:E}, \S\ref{sec:lemmaE}).

The motive is the classical line connecting covering systems to roots of
unity: the Mirsky--Newman--Davenport--Rado theorem --- in an exact covering
system of $\Z$ the largest modulus repeats --- is proved by evaluating a
generating function at roots of unity and reading off pole structure
\cite{MirskyNewman}. Lemmas A and B are the cyclic, three-term analogue, and
the group ring is where that analogy becomes literal.

\subsection{The identity and the translation theorem}\label{sec:gt-identity}

Fix a prime $p$ and a \emph{configuration}: three arithmetic progressions
$A_i=\{a_i+jd_i \bmod p:0\le j<s_i\}$ in $\Z_p$ with $d_i\in\Z_p^{*}$,
$s_i\ge 1$. Write $\mathbf 1_i$ for the indicator of $A_i$ and define the
\emph{signed excess}
\[
\varepsilon(x)\;:=\;\sum_{i=1}^{3}\mathbf 1_i(x)-1\;\in\;\{-1,0,1,2\},
\qquad
E(z)\;:=\;\sum_{x\in\Z_p}\varepsilon(x)\,z^x .
\]
Work in the group ring $R_p:=\Z[z]/(z^p-1)$, and put
\begin{align*}
Q(z)\;&:=\;\sum_{i=1}^{3} z^{a_i}\,\bigl(1-z^{d_i s_i}\bigr)
        \prod_{j\ne i}\bigl(1-z^{d_j}\bigr),\\[\medskipamount]
P(z)\;&:=\;\prod_{i=1}^{3}\bigl(1-z^{d_i}\bigr).
\end{align*}
Both are sparse: $Q$ has at most $8\times 3=24$ terms with coefficients
$\pm1$ (each summand is the inclusion--exclusion of the box
$a_i+\{0,d_is_i\}+\{0,d_j\}+\{0,d_k\}$), and $P$ has at most $8$ terms. Let
$\overline\Phi_p:=\sum_{x\in\Z_p}z^x$ denote the all-ones element of $R_p$
(the reduction of the cyclotomic polynomial $\Phi_p$).

\begin{theorem}[group-ring translation]\label{thm:GT}
\;\textup{(i)} For \emph{every} configuration,
\[
Q \;=\; E\cdot P \qquad\text{in } R_p .
\]
\textup{(ii)} The configuration covers $\Z_p$ if and only if
$\varepsilon\ge 0$ pointwise. Consequently, writing
$\mathrm{ov}:=\sum_i s_i-p$, if $\mathrm{ov}=0$ then covering is equivalent
to $E=0$; if $\mathrm{ov}=1$ to $E=z^e$ for some $e$; and if
$\mathrm{ov}=2$ to $E\in\{z^e+z^f,\;2z^e\}$.
\textup{(iii)} The kernel of multiplication by $P$ on $R_p$ is exactly
$\Z\cdot\overline\Phi_p$, and $E$ is the \emph{unique} $X\in R_p$ with
$Q=XP$ and $X(1)=\mathrm{ov}$.
\end{theorem}

\begin{proof}
(i) Let $\omega=e^{2\pi i/p}$ and evaluate at $z=\omega^\ell$,
$\ell\in\Z_p$. For $\ell\ne 0$, since $d_i\in\Z_p^{*}$,
\[
\widehat{\mathbf 1}_{A_i}(\ell)
=\sum_{j<s_i}\omega^{\ell(a_i+jd_i)}
=\omega^{\ell a_i}\,\frac{1-\omega^{\ell d_i s_i}}{1-\omega^{\ell d_i}},
\]
a geometric sum with nonzero denominator. Summing the function identity
$\sum_i\mathbf 1_i=1+\varepsilon$ over $x\in\Z_p$ against $\omega^{\ell x}$
kills the constant $1$ at $\ell\ne0$, giving
$\sum_i\widehat{\mathbf 1}_i(\ell)=\widehat\varepsilon(\ell)$; multiplying by
$P(\omega^\ell)=\prod_j(1-\omega^{\ell d_j})\ne 0$ yields
$Q(\omega^\ell)=\widehat\varepsilon(\ell)P(\omega^\ell)$, i.e.\ the
polynomial $Q-E\cdot P$ vanishes at every nontrivial $p$-th root of unity.
At $\ell=0$ both sides vanish: $Q(1)=0$ because each summand contains the
factor $1-z^{d_is_i}$ at $z=1$, and $E(1)P(1)=\mathrm{ov}\cdot 0=0$. So
$Q-E\cdot P$ vanishes at all $p$ roots of $z^p-1$; a polynomial with integer
coefficients divisible by $z^p-1$ over $\mathbb{Q}$ (the roots are distinct) is
divisible over $\Z$ (monic divisor, Gauss's lemma), and exponent reduction
modulo $p$ changes a polynomial by multiples of $z^p-1$. Hence
$Q=E\cdot P$ in $R_p$, with no hypothesis beyond $d_i\ne 0$.

(ii) Covering means $\min_x\sum_i\mathbf 1_i(x)\ge 1$, i.e.\ $\varepsilon\ge
0$ pointwise; and $\sum_x\varepsilon(x)=\sum_i s_i-p=\mathrm{ov}$ always.
Pointwise $0\le\varepsilon\le 2$ with total mass $0$, $1$, $2$ leaves exactly
the listed profiles.

(iii) Let $X\in\ker(\text{mult.\ by }P)$. Then $X(\omega^\ell)P
(\omega^\ell)=0$ for all $\ell$, and $P(\omega^\ell)\ne 0$ for $\ell\ne 0$,
so the coefficient vector $(X_x)_{x\in\Z_p}$ satisfies
$\sum_x X_x\omega^{\ell x}=0$ for all $\ell\in\Z_p^{*}$. By character
orthogonality on $\Z_p$, $(X_x)$ lies in the span of the all-ones vector;
integrality gives $X=t\,\overline\Phi_p$ with $t\in\Z$. Conversely
$(1-z^{d})\overline\Phi_p\equiv 0$ in $R_p$ for every $d$ (multiplication by
$z^d$ permutes the all-ones vector), so $\Z\overline\Phi_p\subseteq\ker$.
Finally, if $Q=XP=EP$ then $X-E=t\overline\Phi_p$ and
$X(1)-E(1)=t\,p$; the mass condition $X(1)=\mathrm{ov}=E(1)$ forces $t=0$.
\end{proof}

Part (i) is stronger than the conditional identity one obtains by assuming a
covering: the $24$-term element $Q$ is \emph{always} divisible by the
$8$-term element $P$ in $R_p$, and the quotient is always the signed excess.
Individual summands of $Q$ need not be divisible by $P$ --- the divisibility
is a cancellation among the three boxes, and that cancellation is the
arithmetic content of the covering problem. Part (iii) says the covering
question is a genuine linear-algebra sign question about a uniquely
determined object:
\begin{center}
\emph{the configuration covers $\Z_p$ iff the unique right-mass
solution $X$ of $Q=XP$ has nonnegative coefficients.}
\end{center}

\begin{corollary}\label{cor:signform}
Lemma B of the companion notes is equivalent to: for every normalized
critical configuration at $p=6k+5$ with pairwise $\pm$-distinct differences,
the right-mass solution $E$ of $Q=XP$ has a negative coefficient; Lemma A at
$p=6k+1$ likewise. The lemmas are sign conditions on the quotient
$Q/P\in R_p$.
\end{corollary}

\begin{proposition}[normal form]\label{prop:normalform}
The statement $Q=EP$, the covering property, the sizes, the overlap budget,
and pairwise $\pm$-distinctness are preserved by \textup{(a)} translation
$a_i\mapsto a_i+t$, under which $Q\mapsto z^{t}Q$, $E\mapsto z^{t}E$,
$P\mapsto P$; and \textup{(b)} scaling $(a_i,d_i)\mapsto(\lambda a_i,
\lambda d_i)$, $\lambda\in\Z_p^{*}$, under which all three transform by the
ring automorphism $z\mapsto z^{\lambda}$ of $R_p$. Hence every configuration
is equivalent to one with $a_1=0$, $d_1=1$, and the census below may be run
in that normal form; reflection $z\mapsto z^{-1}$ permutes the solution set.
\end{proposition}

\begin{proof}
Translation: $\varepsilon$ transforms to $x\mapsto\varepsilon(x-t)$, so
$E\mapsto z^tE$, and each summand of $Q$ gains the factor $z^t$. Scaling:
$z\mapsto z^\lambda$ is a ring automorphism because $\lambda$ is coprime to
$p$, and it carries the triple $(z^{a_i},z^{d_is_i},z^{d_j})$ to the scaled
configuration's data since $s_i$ is unchanged and
$\lambda(d_is_i)=(\lambda d_i)s_i$. Covering and $\pm$-distinctness are
properties of the bijections $x\mapsto x+t$ resp.\ $x\mapsto\lambda x$ on
$\Z_p$. Scaling by $d_1^{-1}$ and translating by $-a_1$ produce the normal
form.
\end{proof}

\subsection{Machine verification of the translation}\label{sec:gt-verify}

Theorem~\ref{thm:GT} was verified by an implementation independent of the
prior harness: $Q$, $P$, $E$ are built directly from their definitions as
integer arrays over $\Z_p$ (the $24$ terms by explicit inclusion--exclusion
over the three boxes), the product $E\cdot P$ by circular convolution, and
the covering property by point counts. The census search uses a
bit-mask residual-fit enumeration (for each $(d_2,a_2)$ surviving the
inclusion--exclusion prune, the completing placements $a_3$ are obtained
from the minimal cyclic arc containing the scaled residual), with a full
brute-force cross-check at $p=7$ (all placements, no prune, no fit).

\begin{table}[t]
\centering
\small
\caption{Verification of the translation theorem. ``Identity'' is the check
$Q=E\cdot P$ in $R_p$; ``sharp form'' includes configurations that do not
cover, where $E$ has negative coefficients.}
\label{tab:gtverify}
\begin{tabularx}{\textwidth}{@{}lXXr@{}}
\toprule
check & scope & outcome & count \\
\midrule
identity, Lemma A exceptions & $p=7,13$ (all $52$) & all hold; $E$-profiles
exactly $\{z^e\}$, $\{z^e{+}z^f\}$, $\{2z^e\}$ as classified & 52 \\
identity, census coverings & every covering of \S\ref{sec:gt-census}, all
$15$ primes $5\le p\le 61$ & all hold; $E\ge 0$ & 12{,}256 \\
identity, sharp form & $3{,}000$ random configurations per prime, $p\in\{11,
13,17,19,23,29,31,37\}$, covering or not & all hold (signed $E$) & 24{,}000 \\
sign $\Leftrightarrow$ covering & same sample & agreement & 24{,}000 \\
kernel $=\Z\overline\Phi_p$ & same sample, $X=E\pm\overline\Phi_p$ &
$Q=XP$ holds & 24{,}000 \\
$\pm$-distinct census & $p\in\{11,\dots,59\}\cup\{19,\dots,61\}$, exhaustive
& zero coverings & --- \\
brute-force ground truth & $p=7$, all placements & $272$ normalized
configurations $=$ $68$ set-triples, matching the fast enumerator & --- \\
\bottomrule
\end{tabularx}
\end{table}

Two soundness items from this run are recorded in \S\ref{sec:sound}. First,
the first implementation of the fit test had an off-by-one error in the
minimal-arc length; it was caught immediately by the assertion that every
enumerated covering has $E\ge 0$, and after the fix all counts agreed with
the brute force. Second, the run forced a reconciliation of the prior
failure counts: the ``$50$ at $p=7$ and $2$ at $p=13$'' of the earlier
harnesses are \emph{witness-tuple} counts --- one per (ordered sizes,
$d_2$, $d_3$, $a_2$) with $d_2,d_3$ over positive $\pm$-reps and at least
one completing $a_3$, the shared convention of the earlier search and its
brute-force cross-check --- while the true solution counts are $272$
normalized configurations ($68$ set-triples) at $p=7$ and $8$
($2$ set-triples) at $p=13$, confirmed by the no-prune brute force. The
content of the earlier finding (Lemma A is false at $7$ and $13$; the cells
are closed there by direct verification) is unchanged, and the identity
check ran on the \emph{full} solution set, not on witnesses.

\subsection{The covering census: four families, pinned placements}
\label{sec:gt-census}

The census drops the $\pm$-distinctness hypothesis and enumerates
\emph{all} critical coverings in the normal form $a_1=0$, $d_1=1$: all
ordered size triples from the class-appropriate critical window with
$\sum s_i-p$ in the class-appropriate set (Lemma B: sizes
$\{2k{+}1,2k{+}2\}$, $\mathrm{ov}\in\{0,1\}$; Lemma A: sizes
$\{2k,2k{+}1,2k{+}2\}$, $\mathrm{ov}\in\{0,1,2\}$), over \emph{all}
$(d_2,d_3)\in(\Z_p^{*})^2$. Write $\|d\|:=\min(d,-d \bmod p)$ for the
$\pm$-representative and $q:=(p-1)/2$.

\begin{proposition}[census record, machine]\label{prop:census}
At every prime $p\equiv5\pmod 6$ with $11\le p\le 59$, and at $p=5$, the
census returns exactly $144$ normalized coverings, distributed by the
$\pm$-rep pair $(\|d_2\|,\|d_3\|)$ as
\[
(1,1):48,\qquad (1,q):32,\qquad (2,2):32,\qquad (q,1):32,
\]
independently of $p$. At every prime $p\equiv1\pmod 6$ with $19\le p\le 61$
it returns exactly $1368$, distributed as
\[
(1,1):504,\qquad (1,q):288,\qquad (2,2):288,\qquad (q,1):288,
\]
again independently of $p$. No other family occurs at any of these primes,
and no $\pm$-distinct covering occurs ($0$ at every prime in range,
re-deriving Lemmas A and B at those primes by an independent code path).
At the boundary, $p=5$ already shows the four families (with $q=2$); $p=7$
shows the $k=1$ collapse (nine families, including $272$ $\pm$-distinct
configurations --- the $68$ set-triples of the Lemma A failure); $p=13$
shows the four canonical families with the same counts as at general $p$,
plus small exceptional families (among them the $8$ $\pm$-distinct
instances of the two Lemma A failures), totaling $1520$.
\end{proposition}

\begin{table}[t]
\centering
\small
\caption{Per-multiset family counts, constant in $p$ (Lemma A side, $19\le
p\le 61$; the Lemma B side gives $24/16/16/16=72$ for each of its two
multisets). ``ov'' is $\sum s_i-p$.}
\label{tab:permultiset}
\begin{tabular}{@{}lrrrrrr@{}}
\toprule
multiset & ov & $(1,1)$ & $(1,q)$ & $(2,2)$ & $(q,1)$ & total \\
\midrule
$\{2k,2k,2k{+}1\}$ & $0$ & 24 & 16 & 16 & 16 & 72 \\
$\{2k,2k,2k{+}2\}$ & $1$ & 72 & 32 & 32 & 32 & 168 \\
$\{2k,2k{+}1,2k{+}1\}$ & $1$ & 72 & 48 & 48 & 48 & 216 \\
$\{2k,2k{+}1,2k{+}2\}$ & $2$ & 288 & 160 & 160 & 160 & 768 \\
$\{2k{+}1,2k{+}1,2k{+}1\}$ & $2$ & 48 & 32 & 32 & 32 & 144 \\
\bottomrule
\end{tabular}
\end{table}

Every one of the $12{,}256$ census coverings satisfies the identity of
Theorem~\ref{thm:GT} with $E\ge0$ (Table~\ref{tab:gtverify}). The four
families are structural, not merely counted: they are, in order, the
three-interval family, the parity family, and the two two-block families
(interval plus two-block, with the two-block on either index). Their
placement counts per ordered size triple are also $p$-independent --- the
$(1,1)$ family has $8$ per triple at $\mathrm{ov}=0$ and $24$ at
$\mathrm{ov}=1$ --- and this is the identity-level form of the pinning that
the shear-rigidity argument measured fiberwise. The $p$-independence has a
structural explanation, and it is worth isolating, because it is the
mechanism a proof of Conjecture~\ref{conj:SL} will exploit.

\begin{proposition}[why the counts are $p$-independent]\label{prop:pindep}
Fix a family and a multiset of Table~\ref{tab:permultiset}, and $k\ge1$.
Every census configuration of the parity family $(2,2)$ with at most one
foot per AP has the \emph{integer-parity parametrization}: each of $A_2,
A_3$ meets $I_1=[s_1,p-1]$ in a run of consecutive elements of one integer
parity class of $I_1$, does not wrap, and carries its feet as the one-point
extensions of that run at either end; consequently the configurations are
enumerated by a choice set --- parity assignment ($\le2$), orientation of
each AP ($2$ each), spill slot of each AP ($\le3$ each) --- of cardinality
bounded independently of $k$, and the per-multiset counts are the constants
of Table~\ref{tab:permultiset}. Written out: $16$ for both B-side
multisets, $16$ ($M1$), $48$ ($M2$), $32$ ($M3$), $32$ ($M5$); the $M4$
count $160$ adds the wrapped spill pair $\{0,2\}$ available at
$\mathrm{ov}=2$. The $(1,1)$ family is Proposition~\ref{prop:F1}; the
two-block families obey the same junction arithmetic (verified; the
$s_1=2k{+}1$ slice is written out in \S\ref{sec:pq-census}).
\end{proposition}

\begin{proof}
Let $A=\{a+2j\}$ ($d=\pm2$) have $|A\setminus I_1|\le1$ and meet $I_1$.
The step-$2$ cycle visits $\dots,p{-}3,p{-}1,1,3,\dots$ (or shifted by
one): a wrap passes from $p-1$ (or $p-2$) to $0$ (or $1$). If the segment
wrapped strictly inside, the post-wrap head would contribute at least two
points of $A_1=[0,s_1-1]$ once it re-enters $I_1$ --- impossible at one
foot unless the wrap is at an endpoint of the segment; hence the $I_1$-part
is the pre-wrap tail: consecutive integers of one parity, and the single
foot (if any) is the adjacent run extension below $s_1$ or above $p-1$
(after reduction). Covering forces $A_2\cup A_3\supseteq I_1$, and since
each parity class of $I_1$ can only be covered by the AP of that parity,
one AP carries the even run, the other the odd run, each extended by its
spill slots within the foot budget $\mathrm{ov}$. The count is then a
product of $\{\text{assignment}\}\times\{\text{orientations}\}\times
\{\text{spill slots}\}$, all bounded by constants; evaluating it per
multiset (with $|E(I_1)|$, $|O(I_1)|$ determined by $s_1$ alone) gives the
listed constants, each independent of $k$. E.g.\ the B-side multiset
$\{2k{+}2,2k{+}2,2k{+}1\}$: at $s_1=2k{+}1$ the parities have $2k{+}2$
each, both APs are exact runs ($2$ orientations each, $2$ assignments:
$8$); at $s_1=2k{+}2$ the even run has $2k{+}2$, the odd $2k{+}1$, the
sizes match forced assignments ($2$ ordered triples $\times2^2$
parametrizations: $8$); total $16$. The structural claim was verified on
every census $(2,2)$ solution at six primes ($32$ resp.\ $288$ per prime,
$100\%$ conformity, spills $\le\mathrm{ov}$ in $A_1$).
\end{proof}

This is the answer to ``why $p$-independent'': the critical configurations
live in a finite, $k$-independent choice space of interval-arithmetic
patterns; the families are its connected components, and the placement
counts are its component sizes. The same mechanism governs the row census
of the $pq$ side (\S\ref{sec:pq-census}), where the pinned placement sets
have $24$ resp.\ $312$ elements at every row modulus --- which is why the
covering counts $28/68$ of the shear-rigidity input are $p$-independent
(\S\ref{sec:gt-support}).

\begin{proposition}[the interval family]\label{prop:F1}
In the normal form, with $d_2,d_3\in\{\pm1\}$ and critical sizes with
$\mathrm{ov}\in\{0,1\}$ (for $k\ge1$, so that all sizes are $\ge2$; the
degenerate prime $p=5$ is verified directly), the coverings are exactly the consecutive
arrangements: at $\mathrm{ov}=0$, the two cyclic orders of the consecutive
tiling of $I_1=[s_1,p-1]$ by $A_2,A_3$; at $\mathrm{ov}=1$, the six
arrangements obtained from them by placing the single double point at one of
three positions --- the interior junction, the left spill into $A_1$ at
$s_1-1$, or the wrap at $0$ --- each in two orders. Per ordered size triple
this gives $2$ resp.\ $6$ set-arrangements, each realized by $2^2$
parametrizations (independent orientation choices $\pm1$ for $A_2$ and
$A_3$), i.e.\ $8$ resp.\ $24$ configurations, matching the census exactly.
\end{proposition}

\begin{proof}
With $d_i=\pm1$ each $A_i$ is a cyclic interval. Let $q=p-s_1=|I_1|$. Since
$A_1=[0,s_1-1]$, covering is equivalent to $A_2\cup A_3\supseteq I_1$.
At $\mathrm{ov}=0$ we have $s_2+s_3=q$, so $A_2\cup A_3=I_1$ exactly; both
intervals lie in $I_1$, and two intervals whose union is the arc $I_1$ are
its consecutive blocks, in one of the two orders. At $\mathrm{ov}=1$ we have
$s_2+s_3=q+1$, so $|A_2\cap A_3|=1-(A_2\cup A_3\smallsetminus I_1)$ counts
the single double point (a triple point would contribute $2$ to the mass,
exceeding $\mathrm{ov}=1$). If the double point is interior
($\in I_1$) then $A_2\cup A_3=I_1$ and the two intervals are consecutive
blocks sharing their junction point. If it lies in $A_1$, then
$A_2\cap A_3=\varnothing$, and the one point of $(A_2\cup A_3)\cap A_1$
belongs to exactly one of the intervals: an interval of size $\ge2$
containing a point $<s_1$ and points of $I_1$ must contain the intervening
block of $A_1$, so it either starts at $s_1-1$ (left spill) or wraps through
$0$, in which case it contains $[0,v]$ for some $v$ and $v=0$ is forced.
Each case occurs in the two orders of $(A_2,A_3)$, giving $3\times2=6$
set-arrangements. The parametrization count is immediate.
\end{proof}

The other three families have the same character. The $(2,2)$ family at
$\mathrm{ov}=0$ is the parity split: the surviving structural pair
$(E,O)$ of the Lemma B proof --- the two $d=\pm2$ progressions partition
$I_1$ into its even and odd classes, which is why $\pm$-distinctness is
exactly what kills it. The $(1,q)$ and $(q,1)$ families are the two-block
families: one of $A_2,A_3$ is the two-block set of difference $q$
(the proved two-block structure of the trichotomy: blocks of
$\lceil s/2\rceil$ and $\lfloor s/2\rfloor$ consecutive integers whose right
endpoints differ by exactly $q$), and the partner is an interval. The full
placement windows of these three families are verified but not written out
as propositions; they are the remaining content of the support lemma below.

\subsection{The support lemma and the computational inputs}
\label{sec:gt-support}

The census law suggests the all-$p$ statement that the group-ring programme
targets.

\begin{conjecture}[support lemma: the covering classification]\label{conj:SL}
Let $p=6k+5$, $k\ge0$ (resp.\ $p=6k+1$, $k\ge3$) be prime. Every critical
covering configuration at $p$ in the normal form satisfies
\[
(\|d_2\|,\|d_3\|)\in\{(1,1),\,(1,q),\,(2,2),\,(q,1)\},
\qquad q=(p-1)/2,
\]
and conversely each family is realized at every such prime with the
placement counts of Table~\ref{tab:permultiset}, independent of $p$.
Equivalently: \textup{(i)} no pairwise $\pm$-distinct critical
configuration covers $\Z_p$ (Lemmas B and A), and \textup{(ii)} every
critical covering has a $\pm$-collision, the collision patterns being
exactly the four listed.
\end{conjecture}

Its status, part by part. Clause (i) is \emph{proved} for $p\equiv5\pmod6$,
$p\ge17$ (Lemma B; see below), and for $p\equiv1\pmod6$, $p\ge31$
(Lemma A, unconditionally: its foot-count input is discharged by
Lemma~\ref{lem:E} of \S\ref{sec:lemmaE}); verified to $p\le251$ (prior
harness) and $p\le419$ (the foot census), and re-derived independently
at every prime in $5\le p\le 61$. The $(1,1)$ family is proved (Proposition~\ref{prop:F1}); the $(2,2)$
family is proved (Proposition~\ref{prop:pindep}); the two-block placement
structure is proved. What is \emph{not} proved is the exhaustiveness of the
four families at general $p$ (clause (ii) and the two-block placement
counts), verified by the fresh census to $p\le61$ on both classes (and at
$5$, $7$, $13$) and by the prior relaxed census to $p\le250$. The
$\varepsilon=2$ counts of the interval family are verified but, unlike
Proposition~\ref{prop:F1}, not written out.

\medskip
\noindent\textbf{The computational inputs, precisely.} The $p^2$
closure --- the $(1K,3U)$ cell is empty at every prime $p\ge5$ --- rested,
at the previous revision, on two named inputs beyond the proved lemmas.
Both were \emph{unbounded} hypotheses: each asserts something for all
primes in a mechanism range and is verified only to a finite bound, so
neither can be discharged by a bounded search, however large; each
requires a lemma, and each had one identified. The first of the two
(Input E below) is now \emph{proved} --- Lemma~\ref{lem:E} of
\S\ref{sec:lemmaE} --- leaving Input S as the closure's single
unbounded input.

\emph{Input E (the foot-count trichotomy, stated in its consumed form) ---
discharged in this revision.}
The kills M4 and M5 of Lemma A consume exactly the \emph{one-foot} case, at
exactly two size pairs: M4 at $(s_1,s)=(2k{+}2,\,2k{+}1)$, M5 at
$(2k{+}1,\,2k{+}1)$. The consumed statement is therefore: \emph{at
$p=6k+1$, any critical-size AP with at most one point outside
$I_1=[s_1,p-1]$ at these two pairs has $\|d\|\in\{2,(p{-}1)/2\}$} --- at
the M4 pair the viable set is in fact exactly $\{2\}$, at the M5 pair
exactly $\{2,(p{-}1)/2\}$, matching the kills' branching. A complete
census at every prime $p\le419$ gives the exact boundaries: the consumed
form is clean at every prime $p\equiv1\pmod6$ with $p\ge31$ (dirty at $13$
and $19$). \emph{The input is now discharged}: Lemma~\ref{lem:E} of
\S\ref{sec:lemmaE} (the three-distance foot count) proves the consumed
statement outright at every prime $p\ge31$, by an elementary column
decomposition of the dilation $j\mapsto jd\bmod p$ with no group-ring
content, machine-verified against the census ($76$/$76$ viable-set
agreement at $p\le419$). Consequently \emph{Lemma A is proved for
$p\ge31$}, with $p=19$ closed by the finite case analysis and $7,13$ by
direct verification; the all-pairs one-foot version is clean for $p\ge31$
on the $\epsilon=1$ class and $p\ge41$ on the $\epsilon=5$ class; the two-foot
version (clean only from $p\ge61$ resp.\ $47$) is consumed by \emph{no}
current proof --- the earlier ``$p\ge43$'' phrasing bundled the two-foot
statement with the pair-restricted one, and a previous revision separated
them. Lemma B, finally, consumes \emph{no} foot-count input: its kills run
on the contained census (Lemma T, clean for $p\equiv5\pmod6$, $p\ge17$) and
Corollaries T1--T3, so \emph{Lemma B is proved for $p\ge17$} (the earlier
$p\ge23$ reflected a pre-T3 draft of the $\beta$-kill; at $p=17$ the
contained forms at $(s_1,s)=(6,6)$ are exactly $E$ and $T$ with
$|E\cap T|=4>1$, as the proof requires). The earlier discharge sketch ---
a wrap of a middle difference $d\in[3,s_1]$ landing in $A_1$, the
re-entry climb contributing at least $\approx s\cdot s_1/p-2$ points ---
is superseded by the exact column analysis of \S\ref{sec:lemmaE}, which
proves the bound outright rather than heuristically.

\emph{Input S (the pinning lemma of the shear-rigidity argument).}
Statement: in the $\pm$-colliding closure, the per-fiber covering
configurations of a colliding family are pinned --- the relative-offset
difference sets have at most $6$ elements, with covering counts $28$ and
$68$ per collision topology, independent of $p$. Consumer: the
shear-rigidity theorem for $\pm$-colliding units --- hence the closure for
colliding unit triples at $p\ge11$. Status: verified at every prime
$11\le p\le113$; exhaustive ground truth $3.3$ million families at $p\le43$
with zero coverings; mechanism (exact linear per-fiber drift against
pinned configurations) proved in general. Discharge: the collision side of
Conjecture~\ref{conj:SL}. The per-fiber covering configurations are exactly
critical coverings of the fiber's $t$-space ($\cong\Z_p$) whose three
differences carry the colliding units' $\pm$-collision; the support lemma
classifies these into the four families with $p$-independent placement
sets, and the pinning tables are then a read-off of that classification ---
which is why the counts are $p$-independent: the families are. The $pq$
side measures the pinning directly and confirms it: every row configuration
satisfying the fiber condition is a census solution after normalization
(Proposition~\ref{prop:pqpinned}), with at most $6$ surviving rows per
family against fiber arcs of length $\tfrac23 q$.

In this form the technical debt is finite and explicit: the prime-side
foot-count lemma (Input E, one-foot at two size pairs) is now \emph{proved}
(Lemma~\ref{lem:E}, \S\ref{sec:lemmaE}); what remains is the support
lemma's exhaustiveness clause and the shear input --- each with a
statement, a consumer, a verification range, and a proof strategy; the
group-ring translation is the machine that proves them, and the $pq$ side
of \S\ref{sec:pq} adds no further input. After this revision the entire
programme rests on a single unbounded input, Input S.

\medskip
Finally, the defect statistics. Define the distance-to-solvability of a
non-covering $\pm$-distinct critical configuration as $\|Q\|_{0}$ (the
number of nonzero coefficients of $Q$ in $R_p$) at $\mathrm{ov}=0$, and
$\min_e\|Q-z^eP\|_{0}$ at $\mathrm{ov}=1$ --- the least $\ell_0$-distance
from the identity's data to a covering profile. Over $3{,}000$ random
configurations per prime the sampled minima grow: $5,5,6,7,9$ at
$p=11,17,23,29,41$ (medians $8,11,13,14,17$), with the minimizers
concentrated in the overlap-$1$ multiset $\{2k{+}2\}^3$ --- the case of the
Lemma B proof that needed Corollaries T2/T3 rather than pure parity. The
obstruction thus \emph{strengthens} with $p$: this is the cyclotomic
analogue of the shear-rigidity near-miss margin (families failing all but
two fibers), and it is the quantity a stability argument would bound below.

\subsection{Outlook: cyclotomic tools and the $pq$ side}\label{sec:gt-outlook}

The translation puts the problem in the room of the
Lam--Leung--de~Bruijn--Schoenberg theory of polynomials with few terms
divisible by cyclotomic polynomials
\cite{deBruijnSchoenberg,LamLeung}. The prime-modulus case of that theory is
elementary and worth recording, since it is the tool any proof of
Conjecture~\ref{conj:SL} will consume on the nonnegative side:

\begin{lemma}[prime-modulus Lam--Leung structure]\label{lem:LL}
Let $F\in\Z[z]$ have nonnegative coefficients and satisfy $\Phi_p\mid F$.
Then, reducing exponents modulo $p$, all $p$ residue classes carry equal
mass: $F\equiv(F(1)/p)\,\overline\Phi_p$ in $R_p$.
\end{lemma}

\begin{proof}
$F=\Phi_pG$ gives $p\mid F(1)$, and $F(\omega^\ell)=0$ for $\ell\ne0$ says
the reduced coefficient vector is annihilated by all nontrivial characters
of $\Z_p$; by orthogonality it is constant.
\end{proof}

The gap between this and our problem is exactly the sign structure: $Q$ is
a $\pm1$ combination, not a nonnegative one, so Lemma~\ref{lem:LL} does not
apply directly, and the covering question is a \emph{sign-mixed} sparse
divisibility problem --- the quotient $E=Q/P$ must be shown to have a
negative coefficient in the $\pm$-distinct regime, and to be nonnegative
with the exact four-family profile in the colliding regime. The specific
structure available beyond general sparsity is that $Q$ is built from
factors $(1-z^{c})$: these are real-rooted at $z=1$, and their coefficient
sequences are log-concave; whether that structure propagates through the
alternating combination $Q$ in the way log-concavity propagates through
products --- the mechanism by which stability methods turn counting
statements into theorems \cite{BrandenHuh} --- is the precise question the
stability route must answer. Conjecture~\ref{conj:SL} is calibrated for it:
the four families are its tightness examples, and the defect margins of
\S\ref{sec:gt-support} are its stability margins.

At composite modulus the translation persists with new structure: at
$N=pq$ the group ring is $\Z[\Z_p\times\Z_q]$, the tensor product of the
prime-modulus rings, the identity tensors over the CRT components, and the
kernel of Lemma~\ref{lem:LL} enlarges --- the cyclotomic field of $\Phi_{pq}$
replaces $\mathbb{Q}(\omega_p)$, and with it the Lam--Leung theory acquires its
genuine (nontrivial) content \cite{LamLeung}. This is not a black-box
transfer of the prime side, and it is now carried out: \S\ref{sec:pq}
proves the row translation (Theorem~\ref{thm:R}), verifies the two-unit
and three-unit cells at ten semiprime moduli, and reduces their closure to
the same input structure, now a single unbounded input after the discharge
of Input E --- the tensor structure is the reason the $pq$ side
inherits, rather than re-derives, the prime-side obstruction theory.
